\section*{5. Noether 定理的例外情形}

本节继续在 $\mathbb C$ 上工作, 使用整系数上同调.

\subsubsection*{5.1.}
设 $2n$, $d$ 及 $\underline a=(a_1,\ldots,a_d)$ 满足 $d\geq1$, $a_i\geq2$, 并且维数 $2n$、多重次数 $\underline a$ 的光滑完全交 $V$ 的中间上同调纯为 $(n,n)$ 型. 换言之, 或者 $n=0$, 或者 $(2n;\underline a)$ 为 $(2n;2)$, $(2n;2,2)$ 或 $(2;3)$ 之一（XI 2.9）.

假设 $V$ 是由维数 $2n+1$、多重次数 $(a_1,\ldots,a_{d-1})$（分别 $(a_1,\ldots,a_d)$）的完全交 $X\subset\mathbb P^{2n+d}$（分别 $X\subset\mathbb P^{2n+d+1}$）, 取某个 Lefschetz 线束 $\ell$ 中次数为 $a_d$ 的超曲面截面（分别超平面截面）所得的.

由 (4.3.1), $H^{2n}(V)$ 的消灭部分 $\operatorname{Ev}(V)$ 等于其本原部分, 它由消灭圈在 $\mathbb Z$ 上生成. 它是超平面截面的上同调类 $\eta$ 的幂 $\eta^n$ 的正交补.

因为 $H^{2n}(V)$ 上的相交型是幺模的, 且 $\eta^n$\footnote{原印本在本原性断言中写作 $\eta^{2n}$, 但此类位于次数 $4n$; 上下文的讨论与下一显示式均要求 $\eta^n\in H^{2n}(V)$.} 不等于另一个上同调类 $y$ 的 $b y$（其中 $b>1$）（XI 1.6(iv)）, 所以 $\operatorname{Ev}(V)$ 上相交型的判别式为
\[
 (\eta^n,\eta^n)=\prod_i a_i.
\]

设 $R$ 为 $H^{2n}(V)$ 中由消灭圈及其共轭组成的根系. 关于 $R$ 已知如下资料:
\begin{enumerate}[label=(\alph*)]
\item $R$ 为空, 或为 $A$, $D$ 或 $E$ 型;
\item $R$ 的秩由 XI 1.10 给出;
\item 其判别式为 $\prod_i a_i$. 在表中, 此判别式是对应的单连通复单群的中心的阶. 计算判别式的论证说明, 这个中心是权格对根格的商的对偶, 且为循环群.
\end{enumerate}

若 $n\ne0$, 在每种情形中这些资料都完全决定 $R$. 得到:

\subsubsection*{命题 5.2.}
\begin{enumerate}[label=(\roman*)]
\item 若 $V$ 为偶维二次超曲面, 则 $R$ 为 $A_1$ 型.
\item 若 $V$ 为三次曲面, 则 $R$ 为 $E_6$ 型.
\item 若 $V$ 的维数为 $2n$, 且为两个二次超曲面之交, 则 $R$ 为 $D_{2n+3}$ 型.
\end{enumerate}

\subsubsection*{5.3.}
假设 $n\ne0$. $\operatorname{Ev}(V)$ 中长度为 $\pm2$ 的向量所成的集合 $R^1$ 也是根系, 且具有与 $R$ 相同的性质 (a), (b), (c). 因此 $R=R^1$. 设 $\pi_1$ 为所有所考虑类型的完全交 $V\subset\mathbb P^{2n+d}$ 的参数空间的基本群. $\pi_1$ 在 $\operatorname{GL}(\operatorname{Ev}(V))$ 中的像为 $R$ 的 Weyl 群. 在每种情形中, 它都是三元组
\[
 \bigl(H^{2n}(V),\eta^n,\text{相交型}\bigr).
\]
的自同构群.

\subsubsection*{5.4.}
将 $R$ 换成 5.3 中的 $R^1$ 后, 断言 5.2(ii) 是熟知的 [1]; 对较小的 $n$, 5.2(iii) 至少为 Manin 所知（$n=0$ 显然; 两个二次超曲面之交所成的曲面是 del Pezzo 曲面）.\footnote{原印本此处指向 4.2(ii) 和 4.2(iii), 但所引用的断言是紧前的命题 5.2 的各个情形.} M. Reid [4] 详细研究了两个二次超曲面之交. 他研究了两个二次超曲面之交 $V\subset\mathbb P^{2n+2}$ 上的 $n$ 维线性子空间的几何, 并由此推出 5.2(iii).

\section*{参考文献}

\begin{enumerate}[label={[\arabic*]}]
\item E. Cartan, ``Quelques remarques sur les 28 bitangentes d'une quartique
  plane et les 27 droites d'une surface cubique,'' \emph{Bull. Sci. Math.}
  \textbf{70} (1946), 42--45; \emph{Oeuvres}, 第 I 部, 第 2 卷, 第~1353 页.
\item P. Deligne, ``Th\'eorie de Hodge II,'' \emph{Publ. Math. IHES}
  \textbf{40} (1971), 5--58.
\item S. Lefschetz, \emph{L'analysis situs et la g\'eom\'etrie alg\'ebrique},
  Gauthier--Villars, 1924（重印于 \emph{Selected Papers}, Chelsea）.
\item M. Reid, 博士学位论文, Cambridge, 1972.
\item J.-P. Serre, ``Faisceaux alg\'ebriques coh\'erents,''
  \emph{Ann. of Math.} \textbf{61} (1955), 197--278（引用时简称 FAC）.
\end{enumerate}
